\documentclass[10pt,a4paper]{amsart}
\usepackage[margin=19mm]{geometry}
\usepackage{amsmath,amssymb,mathtools}
\usepackage[scr=rsfs,cal=euler]{mathalfa}
\usepackage{xfrac}
\usepackage[unicode,hidelinks,bookmarks=false]{hyperref}
\newcommand{\ie}{i.e.\ }
\newcommand{\cf}{cf.\ }
\newcommand{\resp}{respectively\ }
\newcommand{\Subsection}{\subsection}
\providecommand{\nobreakdash}{\nobreak\hbox{-}}
\setlength{\emergencystretch}{2em}
\multlinegap0pt
\usepackage{bm}
\usepackage{bbm}
\usepackage{dsfont}
\usepackage{xspace}
\usepackage{cancel}
\usepackage{mathtools}
\usepackage{amsthm}
\makeatletter
\@ifundefined{thm}{\newtheorem{thm}{Theorem}}{}
\makeatother
\title[Energy approximation for some double phase functionals]{Energy approximation for some double phase functionals}
\author{Menita Carozza, Filomena De Filippis, Raffaella Giova, Francesco Leonetti}
\date{Source: arXiv:2501.06499v2 (CC BY 4.0)}
\begin{document}
\maketitle

\noindent\emph{Source and licence.} Energy approximation for some double phase functionals. Version arXiv:2501.06499v2 (CC BY 4.0), distributed under the Creative Commons Attribution 4.0 International licence, as recorded in the arXiv licence metadata for this exact version and shown on the abstract page https://arxiv.org/abs/2501.06499v2. Full rights notice: licenses/arxiv-2501.06499-CC-BY-4.0.txt.\par\smallskip

\noindent\textit{Editorial scope.} Counted unit: the Thm whose source label is \texttt{absence\_of\_Lavrentiev} in the article (source0.tex lines 1211--1221, source label \texttt{absence\_of\_Lavrentiev}), delivered together with the article's own complete original proof of it (source0.tex lines 1224--1273, one \texttt{proof} environment of 50 lines). No mathematics is written, repaired or re-derived: statement and proof are the article's own text line for line. Required context retained uncounted: strict\_convexity (lines 1185--1188); p\_growth\_from\_below (lines 1190--1193); minimality (lines 1195--1198); energy\_approximation\_sec4 (lines 1200--1203); no\_Lavrentiev\_sec4 (lines 1205--1209). Every definition, display and label of those lines is retained unchanged. Omitted from this package and not redistributed: 1--1184, 1189--1189, 1194--1194, 1199--1199, 1204--1204, 1210--1210, 1222--1223, 1274--1460. Nothing inside a retained excerpt is omitted or altered: every retained body is delivered exactly as the article's lines are written, so no omission receipt of any kind accompanies this package. Citations: the retained excerpts carry no citation command at all, so no bibliography entry of the article is transcribed and no reference is redistributed. Reference adaptation: none was needed and none was made; every reference inside the delivered unit resolves inside it, so no unresolved reference or citation is produced. Printed slips kept literal: no printed word, symbol, hypothesis or number of the counted statement or of its proof was altered. Layout and class: the article's own class is \texttt{article} with packages including amsmath, amsthm, amssymb, amscd, epsfig, esint, mathrsfs, graphics, color, wrapfig, verbatim, babel, amsfonts, dsfont; the wrapper is the lane's pinned \texttt{amsart} editorial wrapper at 10pt on A4, which restates the theorem-like environments the retained text needs and adds the article's own macro definitions that the retained text needs (none), with extra wrapper packages (bm, bbm, dsfont, xspace, cancel, mathtools). The delivered numbering is the unit's own. Assets: no retained span contains a figure environment, a graphics-inclusion command, a PDF-page inclusion, a table or a TikZ picture; no image file is copied into this package, the package declares no asset dependency and no image file is redistributed with it. Licence: version arXiv:2501.06499v2 of this article is distributed under the Creative Commons Attribution 4.0 International licence, as recorded in the arXiv licence metadata for this exact version and shown on the abstract page https://arxiv.org/abs/2501.06499v2. A full rights notice is delivered at licenses/arxiv-2501.06499-CC-BY-4.0.txt. Characters: every retained excerpt is pure ASCII, so the delivered engine, XeLaTeX with its default Latin Modern text font, reports no missing character.
 \begin{equation}
     \label{strict_convexity}
     z \to f(x,z) \quad \text{ is strictly convex.}
 \end{equation}

 \begin{equation}
     \label{p_growth_from_below}
    \tilde{c} (|z|^p - 1) \leq f(x,z).
 \end{equation}

\begin{equation}
     \label{minimality}
    \mathcal{F}(u;B) = \inf_{v \in u + W^{1,p}_0(B;\mathbb{R}^N)} \mathcal{F}(v;B).
 \end{equation}

\begin{equation}
    \label{energy_approximation_sec4}
    \exists \{u_k\} \subset Y: \quad u_k \to u \text{ weakly in } W^{1,p}(B;\mathbb{R}^N) \quad \text{ and } \quad \mathcal{F}(u_k;B) \to \mathcal{F}(u;B).
\end{equation}

\begin{equation}
    \label{no_Lavrentiev_sec4}
    \inf_{v \in u + W^{1,p}_0(B;\mathbb{R}^N)} \mathcal{F}(v;B) =
    \inf_{v \in (u + W^{1,p}_0(B;\mathbb{R}^N)) \cap Y} \mathcal{F}(v;B).
\end{equation}

\begin{thm}
\label{absence_of_Lavrentiev} 
Under assumptions \eqref{strict_convexity}, \eqref{p_growth_from_below} and \eqref{minimality}, we have
\begin{equation}
    \label{no_Lav_implies_en_approx}
    \eqref{no_Lavrentiev_sec4}
    \qquad
    \Longrightarrow
    \qquad\eqref{energy_approximation_sec4}
\end{equation}
\end{thm}

\begin{proof}

\noindent
We use minimality \eqref{minimality} and absence of Lavrentiev phenomenon \eqref{no_Lavrentiev_sec4} in order to get 
$$
\mathcal{F}(u;B) = \inf_{v \in (u + W^{1,p}_0(B;\mathbb{R}^N)) \cap Y} \mathcal{F}(v;B).
$$
Then, there exists a sequence  $\{v_k\} \subset (u + W^{1,p}_0(B;\mathbb{R}^N)) \cap Y$ such that 
\begin{equation}
\label{convergence_sec4}
\mathcal{F}(v_k;B) \to \mathcal{F}(u;B).
\end{equation}
Now we need to prove the convergence of $v_k$ to $u$. To this aim, we use  
convergence \eqref{convergence_sec4} and we get $k_1$ such that
$$
 \mathcal{F}(v_k;B) \leq  \mathcal{F}(u;B) + 1,
$$
for every $k \geq k_1$. On the other hand, 
$p$ coercivity \eqref{p_growth_from_below} implies
$$
\tilde{c} \int_B (|Dv_k|^p - 1) dx \leq \mathcal{F}(v_k;B).
$$
These two inequalities merge into
$$
\tilde{c} \int_B (|Dv_k|^p - 1) dx \leq \mathcal{F}(u;B) + 1.
$$
Such an inequality says that $\{v_k - u \}$ is bounded in $W^{1,p}_0(B;\mathbb{R}^N))$; since $p>1$, there exists a limit function $w_0 \in W^{1,p}_0(B;\mathbb{R}^N))$ and a subsequence 
$\{v_{h_k} - u \}$ such that
$$
v_{h_k} - u \to w_0 \text{ weakly in } W^{1,p}_0(B;\mathbb{R}^N),
$$
\begin{equation}
\label{weak_conv}
v_{h_k}  \to u + w_0 \text{ weakly in } W^{1,p}(B;\mathbb{R}^N),
\end{equation}
\begin{equation}
\label{weak_conv_of_gradients}
Dv_{h_k}  \to D(u + w_0) \text{ weakly in } L^{p}(B;\mathbb{R}^N).
\end{equation}
Convexity \eqref{strict_convexity} implies lower semicontinuity with respect to weak convergence of gradients, so, \eqref{weak_conv_of_gradients} and \eqref{convergence_sec4} give
$$
\mathcal{F}(u + w_0;B) \leq \liminf \mathcal{F}(v_{h_k};B) 
= \lim \mathcal{F}(v_{k};B) = \mathcal{F}(u;B).
$$
Since $u$ is a minimizer and $u + w_0 \in u + W^{1,p}_0(B;\mathbb{R}^N)$, then $u + w_0$ is a minimizer too. Strict convexity \eqref{strict_convexity} implies uniqueness of minimizer, so, $w_0 = 0$ and weak convergence \eqref{weak_conv} tells us that
$$
v_{h_k}  \to u  \text{ weakly in } W^{1,p}(B;\mathbb{R}^N).
$$
We take $u_k = v_{h_k}$: this ends the proof.
\end{proof}

\end{document}
